%%% FIELD proof-local-chart-membership
Put \(c=(a,x)\). Because \(L\ge2\), the support contains two distinct values. Since \(q^0_{c\ell}\ge\eta>0\) by \(\mathrm{ass:eta\text{-}interior}\), every support point has positive probability and hence \(\sigma_c^2(q^0)>0\). Define
\[
m_{c\ell}=\frac{q^0_{c\ell}\{y_\ell-\mu_c(q^0)\}}{\sigma_c^2(q^0)}.
\tag{1}
\]
The definitions of \(\mu_c(q^0)\) and \(\sigma_c^2(q^0)\) give
\[
\sum_{\ell=1}^L m_{c\ell}=0,
\qquad
\sum_{\ell=1}^L y_\ell m_{c\ell}=1.
\tag{2}
\]
Thus the cellwise map from the simplex tangent to its mean is onto. Because \(\operatorname{rank}(F)=d\) by \(\mathrm{ass:menu\text{-}rank}\), the matrix \(A_{\mathrm{map}}=[-F\;F]\) has row rank \(d\). Since \(B_{\mathrm{mean}}\) is onto by (1)--(2), \(A_{\mathrm{tan}}=A_{\mathrm{map}}B_{\mathrm{mean}}\) is onto as well.

Fix \(C_v<\infty\). The set \(\mathcal H\times\{v:\lVert v\rVert\le C_v\}\) is compact, and \((h,v)\mapsto S_0h+Wv\) is linear and continuous. Consequently
\[
C_s(C_v):=\sup_{h\in\mathcal H,\,\lVert v\rVert\le C_v}
\lVert S_0h+Wv\rVert_\infty<\infty.
\tag{3}
\]
If \(n>\{C_s(C_v)/\delta_q\}^2\), then, for every \((a,x,\ell)\),
\[
q_{n,ax\ell}(h,v)
=q^0_{ax\ell}+\frac{s(h,v)_{ax\ell}}{\sqrt n}
\ge q^0_{ax\ell}-\frac{C_s(C_v)}{\sqrt n}
>q^0_{ax\ell}-\delta_q\ge\eta.
\tag{4}
\]
Here \(\delta_q>0\) is \(\mathrm{ass:strict\text{-}interior\text{-}slack}\). Moreover \(s(h,v)\in\mathcal E\), so every cell of \(q_n(h,v)\) sums to one. The three member properties in \(\mathrm{def:saturated\text{-}model}\), together with (4), therefore give \(P_{p,q_n(h,v)}\in\mathcal M_\eta(p)\), uniformly on the displayed compact set.

Linearity of the cell means and the definitions of \(S_0\) and \(W\) give
\[
\sqrt n\{\theta_p(q_n(h,v))-\theta_p(q^0)\}
=A_{\mathrm{tan}}s(h,v)
=A_{\mathrm{tan}}S_0h+A_{\mathrm{tan}}Wv=h.
\tag{5}
\]
By \(\mathrm{ass:tied\text{-}baseline}\), \(\theta_p(q^0)=0\), so (5) proves \(\sqrt n\,\theta_p(q_n(h,v))=h\).

It remains to verify the advertised nuisance coverage. For arbitrary \(h\in\mathcal H\) and bounded \(u\in\mathbb R^{K-d}\), prescribe the cell-mean tangent blocks \(b_0=0\) and \(b_1=Qh+Nu\). Applying (1) cellwise supplies \(s_b\in\mathcal E\) with \(B_{\mathrm{mean}}s_b=b\). Since \(FQ=I_d\) and \(FN=0\),
\[
A_{\mathrm{tan}}s_b=A_{\mathrm{map}}b=F(Qh+Nu)=h.
\]
Thus \(s_b-S_0h\in\ker(A_{\mathrm{tan}})=\operatorname{im}(W)\), and bijectivity of \(W\) gives a unique \(v=W^{-1}(s_b-S_0h)\). Boundedness of this \(v\) on bounded \((h,u)\)-sets follows from continuity of the inverse linear map in finite dimension.

More generally, a contrast-preserving mean tangent satisfies \(F(b_1-b_0)=0\). It decomposes uniquely as a common-arm-mean direction \(b_0=b_1=c\) plus a treatment-effect direction \(b_0=0\), \(b_1=Nu\), because \(\operatorname{im}(N)=\ker(F)\) and \(N\) has full column rank. After fixing these cell means, the remaining tangent in each cell is exactly a vector \(k_c\) satisfying
\[
\sum_{\ell=1}^L k_{c\ell}=\sum_{\ell=1}^L y_\ell k_{c\ell}=0.
\tag{6}
\]
The two linear functionals in (6) are independent because the support contains two distinct values, so this cellwise zero-mean shape space has dimension \(L-2\). In particular, it is the zero-dimensional space when \(L=2\). Across all cells, the common-arm-mean, treatment-effect-null, and shape subspaces have dimensions \(K\), \(K-d\), and \(2K(L-2)\), respectively. Their sum is
\[
K+(K-d)+2K(L-2)=2K(L-1)-d=\dim\ker(A_{\mathrm{tan}}),
\]
where the last equality uses surjectivity of \(A_{\mathrm{tan}}\). The explicit decomposition above shows that their sum is direct and exhausts the kernel. Hence every common-arm-mean and zero-mean outcome-shape direction, including the vacuous shape component for binary outcomes, is represented by the nuisance coordinate \(v\), as claimed.
%%% FIELD proof-coupled-measurable-attainment
Fix \(C_v<\infty\), and write \(c_{ax}=p_xe_{ax}\). The stipulated spaces of \(p\) and \(e\) give \(c_{ax}>0\) for every one of the finitely many cells. Choose two distinct support values \(y_u\ne y_w\), which exist because \(L\ge2\) and the elements of \(\mathcal Y\) are distinct. For any probability vector \(q_{ax}\), direct expansion of the variance gives
\[
\sigma^2_{ax}(q)=\sum_{\ell<m}q_{ax\ell}q_{axm}(y_\ell-y_m)^2.
\tag{7}
\]
Thus \(q^0_{ax\ell}\ge\eta\), by \(\mathrm{ass:eta\text{-}interior}\), implies
\[
\underline\sigma^2:=\eta^2(y_u-y_w)^2
\le \sigma^2_{ax}(q^0)\le 4C_Y^2.
\tag{8}
\]
The upper inequality follows from \(|Y-EY|\le2C_Y\). For \(z\in\mathbb R^d\), the definition of \(\Sigma(r)\) yields
\[
z^\top\Sigma(r)z
=\sum_{x=1}^K\sum_{a=0}^1
\frac{p_x\sigma^2_{ax}(q^0)}{e_{ax}r_{ax}}(z^\top d_x)^2.
\tag{9}
\]
Because \(r_{ax}\in[\rho,1]\), all coefficients in (9) have fixed positive lower and finite upper bounds. Moreover, \(\operatorname{rank}(F)=d\) by \(\mathrm{ass:menu\text{-}rank}\), and multiplication by the positive diagonal matrix \(\operatorname{diag}(p)\) does not change row rank, so \(d_1,\ldots,d_K\) span \(\mathbb R^d\). The continuous function \(z\mapsto\sum_x(z^\top d_x)^2\) therefore has a positive minimum on the Euclidean unit sphere. Equations (8)--(9) consequently place \(\{\Sigma(r):r\in\mathcal R_B\}\) in a fixed compact subset \(\mathcal K\subset\mathbb S_{++}^d\). This argument uses only two distinct support values, so it is unchanged for \(L=2\).

We next construct the pilot objective and prove its uniform accuracy under \(\mathrm{ass:pilot\text{-}rate}\). Put
\[
b_n=m_n^{-1/4}+n^{-1/2},\qquad t_n=b_n^{1/2}.
\tag{10}
\]
Since \(m_n\to\infty\),
\[
b_n\to0,\qquad t_n\to0,\qquad \frac{b_n}{t_n}=b_n^{1/2}\to0.
\tag{11}
\]
Let \(N_{ax}\) be the number of fully labelled pilot records in cell \((a,x)\), and let \(\widehat\nu_{s,ax}\) be the within-cell average of \(Y^s\), for \(s=1,2\), with an arbitrary value in \([-C_Y^s,C_Y^s]\) when \(N_{ax}=0\). If \(\nu_{s,ax,n}=E_{q_n(h,v)}(Y^s\mid A=a,X=x)\), then independence and boundedness give
\[
\operatorname{Var}\!\left(\frac{N_{ax}}{m_n}\right)\le\frac{c_{ax}}{m_n},
\qquad
E\!\left[\left\{\frac1{m_n}\sum_{i=1}^{m_n}
\mathbf1\{X_i=x,A_i=a\}(Y_i^s-\nu_{s,ax,n})\right\}^2\right]
\le\frac{c_{ax}(2C_Y^s)^2}{m_n}.
\tag{12}
\]
Chebyshev's inequality applied first to \(N_{ax}/m_n\) gives \(P(N_{ax}<m_nc_{ax}/2)=O(m_n^{-1})\). On its complement, applying Chebyshev at threshold proportional to \(m_n^{-1/4}\) to the centered numerator in (12), and dividing by \(N_{ax}/m_n\ge c_{ax}/2\), gives
\[
\sup_{h\in\mathcal H,\,\lVert v\rVert\le C_v}
P_{q_n(h,v)}\!\left(
|\widehat\nu_{s,ax}-\nu_{s,ax,n}|>C_{ax,s}m_n^{-1/4}
\right)=O(m_n^{-1/2}).
\tag{13}
\]
The constants do not depend on \((h,v)\), because only the fixed bounds on \(Y\) and the fixed cell probabilities \(c_{ax}\) were used. A union bound over the fixed \(4K\) moment coordinates gives events \(E_n(C_v)\) such that
\[
\inf_{h\in\mathcal H,\,\lVert v\rVert\le C_v}
P_{q_n(h,v)}\{E_n(C_v)\}\to1.
\tag{14}
\]

Linearity of \(s(h,v)=S_0h+Wv\), compactness of \(\mathcal H\), and boundedness of \(\{v:\lVert v\rVert\le C_v\}\) imply a finite \(D(C_v)\) with \(\sup_{h,v}\lVert s(h,v)\rVert_\infty\le D(C_v)\). Since \(q_n-q^0=s(h,v)/\sqrt n\) and the support is finite,
\[
\max_{a,x,s\in\{1,2\}}|\nu_{s,ax,n}-\nu_{s,ax,0}|
\le C(C_v)n^{-1/2}.
\tag{15}
\]
The variance map \((\nu_1,\nu_2)\mapsto\nu_2-\nu_1^2\) is Lipschitz on the bounded moment rectangle. Therefore, after clipping the pilot variance estimates to \([\underline\sigma^2/2,4C_Y^2]\), (13)--(15) imply on \(E_n(C_v)\)
\[
\max_{a,x}|\widehat\sigma^2_{ax}-\sigma^2_{ax}(q^0)|
\le C_0(C_v)b_n.
\tag{16}
\]
This is the only pilot-moment event needed for design selection, and its proof uses \(m_n\to\infty\), with no rate comparison between \(m_n\) and \(\sqrt n\).

Let \(\widehat\Sigma_n(r)\) be the formula for \(\Sigma(r)\) with \(\widehat\sigma^2_{ax}\) in place of \(\sigma^2_{ax}(q^0)\), and let \(\widehat L_n(r)\) be the corresponding Gaussian minimax value. Since \(r_{ax}\ge\rho\), (16) gives
\[
\sup_{r\in\mathcal R_B}\lVert\widehat\Sigma_n(r)-\Sigma(r)\rVert
\le C_1(C_v)b_n.
\tag{17}
\]
For \(S,S'\) in a sufficiently small fixed compact enlargement of \(\mathcal K\), differentiate the centered normal density along \(S_u=(1-u)S+uS'\). The derivative identity
\[
\frac{d}{du}\phi_{S_u}(z)
=\frac12\phi_{S_u}(z)
\left[z^\top S_u^{-1}(S'-S)S_u^{-1}z
-\operatorname{tr}\{S_u^{-1}(S'-S)\}\right]
\]
and the common positive lower and finite upper eigenvalue bounds imply, after integration in \(z\) and \(u\),
\[
\sup_{h\in\mathbb R^d}
\operatorname{TV}\{N(h,S),N(h,S')\}
\le C_2\lVert S-S'\rVert.
\]
The regret is bounded by \(R_H=2\sup_{h\in\mathcal H}\lVert h\rVert_\infty<\infty\). Hence, for every rule \(\delta\) and state \(h\), changing \(S\) to \(S'\) changes expected regret by at most \(R_HC_2\lVert S-S'\rVert\). Taking the supremum over \(h\), then the infimum over \(\delta\), in both directions proves from (17)
\[
\sup_{r\in\mathcal R_B}|\widehat L_n(r)-L(r)|
\le C_3(C_v)b_n
\quad\text{on }E_n(C_v).
\tag{18}
\]

We now give the measurable nested-cell selector. Set \(k_n=\lfloor\log_2\log(n+2)\rfloor\). Through depth \(k_n\), use exactly the coordinate order, parent-only recursion, closed-child convention, and lexicographic order in \(\mathrm{def:implementation\text{-}handle}\). For each nonempty feasible part of a child of the currently selected parent, approximate its minimum \(\widehat L_n\)-value with deterministic error at most \(b_n\). Such approximations can be constructed by successively refining finite state nets, Gaussian observation rectangles, and the projected design nets supplied by \(P_B\).

To verify the boundary case, continuity of \(P_B\), its identity action on \(\mathcal R_B\), and density of the ambient full-box nets make their projected images dense in \(\mathcal R_B\). If \(C=\mathcal R_B\cap D\) is a nonempty closed child cell, send each projected net point \(g\) to its Euclidean nearest point \(c_C(g)\in C\). Existence follows from compactness; uniqueness follows because \(C\) is convex and squared Euclidean distance is strictly convex. If \(c\in C\) and \(\lVert g-c\rVert\le\xi\), then
\[
\lVert c_C(g)-c\rVert
\le\lVert c_C(g)-g\rVert+\lVert g-c\rVert\le2\xi,
\]
so these cell projections form refining finite nets even when \(C\) is a singleton or lies wholly on a dyadic boundary. For effectively represented instances, \(\mathrm{thm:certified\text{-}computation}\), applied after appending the child's rational box inequalities, gives terminating outward certificates. For a fixed mathematical instance, deterministic simple-function approximations give the same existence conclusion without asserting an effective stopping rule: the Gaussian total-variation bound above gives uniform covariance continuity and compactness of \(\mathcal H\) gives finite state nets. Every finite-stage operation is a Gaussian rectangle integral, a finite sum, a cell projection, or a finite linear program. The child index is finite at each stage; resolving its finitely many ties lexicographically makes the resulting map of the pilot moments Borel measurable.

The game approximation is uniform rather than pointwise. Differentiating a Gaussian density in its mean gives a constant \(C_5\), uniform over the compact covariance family, such that
\[
\operatorname{TV}\{N(h,S),N(h',S)\}\le C_5\lVert h-h'\rVert.
\]
Since \(h\mapsto\ell_{\mathrm{reg}}(h,j)\) is Lipschitz uniformly over \(j\in[M]\), every rule's risk has one common modulus of continuity in \(h\). Thus replacing \(\mathcal H\) by a sufficiently fine finite net changes every worst-state risk by a uniformly vanishing amount. For a finite state net and a finite covariance net, let \(\mu\) be the sum of the finitely many Gaussian laws involved. Rational rectangles generate the Borel sigma-field, so each component of any measurable simplex-valued rule has rectangle-simple approximants converging in \(L^1(\mu)\); projecting their vector of values onto \(\Delta_M\) preserves measurability and changes the finite collection of risks by a quantity tending to zero. A covariance net and the Gaussian total-variation bound extend this approximation uniformly over \(\mathcal K\). Taking a rule within an arbitrary positive amount of the defining infimum of \(L\) proves that the finite rectangle games converge uniformly to the Gaussian games. This justifies both the deterministic cell-value accuracy and the finite linear-program representation.

Let \(\widehat v_{n,j,b}\) be the resulting approximate value of feasible child \(b\) at depth \(j\), and, when the parent is \(\mathcal D_{j-1}^\dagger\), let
\[
v_{j,b}=\min_{r\in\mathcal R_B\cap D_{j,b}}L(r).
\]
Equation (18) and the deterministic numerical tolerance give a deterministic \(\alpha_n=C_4(C_v)b_n\) such that, on \(E_n(C_v)\),
\[
|\widehat v_{n,j,b}-v_{j,b}|\le\alpha_n
\tag{19}
\]
simultaneously for every child tested through depth \(k_n\). By (10)--(11), \(\alpha_n=o(t_n)\). Retain precisely the siblings satisfying
\[
\widehat v_{n,j,b}\le\min_c\widehat v_{n,j,c}+t_n,
\tag{20}
\]
and select the lexicographically first retained child. A sibling attaining the finite approximated minimum is always retained, so all empirical cells are nonempty; the parent-only recursion makes them nested.

For the feasible children of the population parent \(\mathcal D_{j-1}^\dagger\), define
\[
E_j=\{b:v_{j,b}=V^*\},\qquad
\gamma_j(L)=
\begin{cases}
\min_{b\notin E_j}\{v_{j,b}-V^*\},&
E_j\text{ is not the set of all feasible children},\\
+\infty,&E_j\text{ is the set of all feasible children}.
\end{cases}
\tag{21}
\]
At least one child is in \(E_j\): the finitely many closed children cover the selected parent cell, which contains a minimizer by the definitions of \(\mathcal D_{j-1}^\dagger\) and \(\mathcal C_{j-1}^\dagger\). Each feasible child is compact, and continuity of \(L\), which follows from the Gaussian total-variation bound and the continuous covariance formula, makes its displayed minimum attained. Consequently every ineligible child has value strictly above \(V^*\); finiteness of the sibling set gives \(\gamma_j(L)>0\).

Suppose inductively that the empirical parent at depth \(j-1\) is \(\mathcal D_{j-1}^\dagger\). By (19), every child in \(E_j\) lies within \(2\alpha_n\) of the smallest empirical value and is retained if \(2\alpha_n<t_n\). Every child outside \(E_j\) lies at least \(\gamma_j(L)-2\alpha_n\) above the smallest empirical value and is rejected if
\[
t_n+2\alpha_n<\gamma_j(L).
\tag{22}
\]
Under these inequalities the retained set equals \(E_j\), so its lexicographically first element is exactly \(\mathcal D_j^\dagger\). Induction proves stabilization of every fixed prefix. If a population minimizer is on a shared dyadic boundary, every closed child containing it is separately feasible and belongs to \(E_j\); (20) retains all such children before the same lexicographic convention is applied. Disconnected, flat, or tied minimizing sets merely enlarge \(E_j\). Thus the argument includes both separated tied minima and dyadic-boundary minima, establishing the stated early-kill property.

The outward cell bounds in \(\mathrm{thm:certified\text{-}computation}\), or the same deterministic convergent bounds for a fixed instance, eventually give a positive lower certificate for every positive gap in (21). Let \(j_n(L)\) be the largest \(j\le k_n\) for which the certified lower gap bounds exceed \(2t_n\) at every level \(1\le i\le j\), with \(j_n(L)=0\) if none is certified. This definition depends only on the pilot schedule and the population prefix-separation profile, not on \(C_v\). For each fixed \(C_v\), (11) implies that eventually \(2\alpha_n<t_n\), so \(2t_n<\gamma_i(L)\) implies (22). Every finite minimum \(\min_{i\le j}\gamma_i(L)\) is positive, whereas \(k_n\to\infty\) and \(t_n\to0\). Hence
\[
j_n(L)\to\infty,
\qquad
a_n(L):=\sup_{m\ge n}2^{-j_m(L)}\downarrow0.
\tag{23}
\]
On \(E_n(C_v)\), the empirical and population paths agree through depth \(j_n(L)\). Choose the final feasible point of a selected cell by successively minimizing its coordinates over that compact cell; since the cell indices form a countable discrete set, this is a measurable fixed lookup. Both this point \(\widehat r_n\) and \(r^\dagger\) lie in the same depth-\(j_n(L)\) population cell, even if later empirical descendants select another minimizing branch within it. The dyadic diameter bound gives
\[
\lVert\widehat r_n-r^\dagger\rVert
\le\sqrt{2K}(1-\rho)2^{-j_n(L)}
\le\sqrt{2K}(1-\rho)a_n(L).
\tag{24}
\]
Combining (14), (23), and (24) proves locally uniform convergence and \(\lVert\widehat r_n-r^\dagger\rVert=O_P\{a_n(L)\}\). The sequence is necessarily schedule- and law-specific through \(b_n\) and the prefix gaps \(\gamma_j(L)\).

It remains to couple this design selector to the terminal rule. At deterministic refinement level \(s\), use the pilot covariance \(\widehat\Sigma_n(\widehat r_n)\), take a finite state net, a rational observation-rectangle partition after tail truncation, and solve the resulting finite primal--dual Gaussian game to total error at most \(2^{-s}\); resolve all finite ties lexicographically. The construction in \(\mathrm{thm:certified\text{-}computation}\) supplies these objects on represented instances, and the deterministic refinements above supply them for the fixed mathematical instance. Let \(J_s<\infty\) be the number of rectangles. A diagonal sequence \(s_n\uparrow\infty\) can be chosen sufficiently slowly that
\[
J_{s_n}n^{-1/2}\to0,
\tag{25}
\]
while its state-mesh, tail, quadrature, rounding, covariance, and finite-game errors all tend to zero. Because each level uses finitely many Borel pilot statistics and lexicographic comparisons, this rectangle-constant rule is jointly measurable with \(\widehat r_n\).

For the distributional coupling, let \(\beta_n=\beta(q_n(h,v))\), \(b=\beta_n-\widehat\beta\), and, conditional on \(\mathscr P_n\), define
\[
H_{i,ax}=\frac{\mathbf1\{X_i=x,A_i=a,R_i=1\}}
{p_xe_{ax}\widehat r_{n,ax}}
(Y_i^{\mathrm{obs}}-\widehat\beta_{ax}).
\tag{26}
\]
Annotation ignorability gives \(E(H_i\mid\mathscr P_n)=b\). The summands \(A_{\mathrm{map}}(H_i-b)\) are conditionally centered and uniformly bounded, because \(|Y|\le C_Y\), \(\widehat\beta\) is clipped to the convex hull of \(\mathcal Y\), \(p_xe_{ax}>0\), and \(\widehat r_{n,ax}\ge\rho\). Only one coordinate of \(H_i\) can be nonzero, so direct multiplication gives their exact covariance
\[
S_n=A_{\mathrm{map}}\left[
\operatorname{diag}_{a,x}\left\{
\frac{\sigma^2_{ax}(q_n)+b_{ax}^2}
{p_xe_{ax}\widehat r_{n,ax}}
\right\}-bb^\top\right]A_{\mathrm{map}}^\top.
\tag{27}
\]
The last two terms inside the brackets form a covariance matrix and are positive semidefinite. Applying the same second-moment calculation as in (12) to the pilot cell means yields, locally uniformly,
\[
\lVert b\rVert=O_P(m_n^{-1/2}).
\tag{28}
\]
The affine chart gives \(\max_{a,x}|\sigma^2_{ax}(q_n)-\sigma^2_{ax}(q^0)|=O(n^{-1/2})\) uniformly on the bounded fiber. Inserting the columns \(-p_xd_x\) and \(p_xd_x\) of \(A_{\mathrm{map}}=[-F\;F]\) shows exactly that
\[
A_{\mathrm{map}}\operatorname{diag}_{a,x}\left\{
\frac{\sigma^2_{ax}(q^0)}{p_xe_{ax}r_{ax}}
\right\}A_{\mathrm{map}}^\top=\Sigma(r).
\tag{29}
\]
Equations (24), (27)--(29), the annotation floor, and continuity of \(r\mapsto1/r\) on \([\rho,1]\) imply
\[
\lVert S_n-\Sigma(r^\dagger)\rVert
=O_P\{m_n^{-1}+n^{-1/2}+a_n(L)\}=o_P(1)
\tag{30}
\]
locally uniformly. This is where pilot estimation enters the covariance: \(m_n^{-1}\to0\) suffices.

By (8), the uniform \(O(n^{-1/2})\) variance perturbation from \(q^0\) to \(q_n\), and the rank argument following (9), \(S_n\) has a common positive lower eigenvalue with probability tending to one uniformly on the bounded fiber. Normalize the \(n\) summands in
\[
\sqrt n\{\widetilde\theta_n-\theta_p(q_n(h,v))\}
=\frac1{\sqrt n}\sum_{i=1}^nA_{\mathrm{map}}(H_i-b)
\tag{31}
\]
by \(S_n^{-1/2}\). Uniform boundedness makes the sum of their third Euclidean absolute moments \(O(n^{-1/2})\). The frozen result \(\mathrm{lem:bentkus\text{-}convex\text{-}set}\) therefore bounds the conditional probability error of each observation rectangle by \(O(n^{-1/2})\). Summing over the \(J_{s_n}\) rectangles and using (25), then using (30) and the Gaussian total-variation bound, gives
\[
\sup_{h\in\mathcal H,\,\lVert v\rVert\le C_v}
|\operatorname{Risk}_n(h,v)-\operatorname{Risk}_{G,n}(h,v)|=o(1).
\tag{32}
\]
Here \(\operatorname{Risk}_{G,n}(h,v)\) is the risk of the same rectangle rule when its input has law \(N(\sqrt n\,\theta_p(q_n(h,v)),\Sigma(r^\dagger))\); \(\mathrm{thm:local\text{-}chart\text{-}membership}\) identifies this mean with \(h\). Bounded regret turns the locally uniform conditional convergence in probability into unconditional convergence. This proves the jointly refined Gaussian-rule risk coupling. The separate condition \(m_n/n\to0\) in \(\mathrm{ass:pilot\text{-}vanishes}\) keeps the pilot likelihood negligible in \(\mathrm{thm:nuisance\text{-}uniform\text{-}experiment}\); no stronger pilot growth is used.

Finally, the claimed limitation is only generic. For any \(\epsilon>0\), define on \([0,1]\)
\[
f_+(z)=\epsilon z,\qquad f_-(z)=\epsilon(1-z),\qquad g(z)=\epsilon/2.
\tag{33}
\]
The unique canonical minimizers of \(f_+\) and \(f_-\) are zero and one, while
\[
\lVert g-f_+\rVert_\infty=\lVert g-f_-\rVert_\infty=\epsilon/2.
\tag{34}
\]
Any selector seeing only the common approximation \(g\) returns the same \(z_g\) in the two cases, and
\[
\max\{|z_g|,|z_g-1|\}\ge\frac12.
\tag{35}
\]
Letting \(\epsilon\downarrow0\) proves that continuity plus uniform objective-approximation control alone gives no spatial recovery modulus tending to zero uniformly over all continuous objectives. Because (33) is not asserted to lie in the Gaussian-risk family \(r\mapsto L(r)\), this establishes no model-specific impossibility, exactly as claimed.
%%% FIELD proof-feasible-risk-attainment
Use the jointly measurable construction proved in \(\mathrm{thm:coupled\text{-}measurable\text{-}attainment}\): fully label the independent pilot, select \(\widehat r_n\in\mathcal R_B\), annotate each main-wave outcome independently at its selected cell rate, compute \(\widetilde\theta_n\), and apply the coupled rectangle-constant Gaussian rule to \(\sqrt n\,\widetilde\theta_n\). Feasibility holds sample by sample. The condition \(\mathrm{ass:pilot\text{-}rate}\) gives \(m_n\to\infty\), while \(\mathrm{ass:pilot\text{-}vanishes}\) gives \(m_n/n\to0\). The former is enough for the design and nuisance estimates, and the latter makes the pilot likelihood asymptotically negligible in \(\mathrm{thm:nuisance\text{-}uniform\text{-}experiment}\). Together with \(\widehat r_n\to r^\dagger\), these facts put the constructed sequence \(\mathcal A^*\) in \(\mathfrak A_{\mathrm{det}}\). The prescribed coordinate order, parent-only recursion, separate testing of shared-boundary children, stabilization of every fixed prefix, and the rate \(O_P\{a_n(L)\}\) follow from \(\mathrm{thm:coupled\text{-}measurable\text{-}attainment}\).

We verify the estimator limit and its precise pilot requirement. Fix \(C_v<\infty\), and define
\[
c_{ax}=p_xe_{ax},\qquad
\beta_n=\beta(q_n(h,v)),\qquad
b=\beta_n-\widehat\beta,
\]
and, conditionally on \(\mathscr P_n\),
\[
H_{i,ax}=\frac{\mathbf1\{X_i=x,A_i=a,R_i=1\}}
{c_{ax}\widehat r_{n,ax}}
(Y_i^{\mathrm{obs}}-\widehat\beta_{ax}).
\tag{36}
\]
By \(\mathrm{ass:randomization}\) and \(\mathrm{ass:consistency}\), the conditional observed-outcome mean is \(\beta_{n,ax}\). By \(\mathrm{ass:annotation\text{-}ignorability}\) and the selected Bernoulli acquisition probability,
\[
E(H_{i,ax}\mid\mathscr P_n)
=\frac{c_{ax}\widehat r_{n,ax}}
{c_{ax}\widehat r_{n,ax}}
(\beta_{n,ax}-\widehat\beta_{ax})=b_{ax}.
\tag{37}
\]
Using the definitions of \(\widetilde\beta_n\), \(\widetilde\theta_n\), and \(\theta_p\), (37) yields the exact identities
\[
E(\widetilde\beta_n\mid\mathscr P_n)=\beta_n,
\qquad
\sqrt n\{\widetilde\theta_n-\theta_p(q_n(h,v))\}
=\frac1{\sqrt n}\sum_{i=1}^nA_{\mathrm{map}}(H_i-b).
\tag{38}
\]

Only one coordinate of \(H_i\) is nonzero. Its conditional second moment and (37) therefore give the exact covariance
\[
S_n=A_{\mathrm{map}}\left[
\operatorname{diag}_{a,x}\left\{
\frac{\sigma^2_{ax}(q_n(h,v))+b_{ax}^2}
{p_xe_{ax}\widehat r_{n,ax}}
\right\}-bb^\top\right]A_{\mathrm{map}}^\top.
\tag{39}
\]
The bracketed remainder
\[
\operatorname{diag}_{a,x}\left\{
\frac{b_{ax}^2}{p_xe_{ax}\widehat r_{n,ax}}
\right\}-bb^\top
\tag{40}
\]
is positive semidefinite because it is the covariance of the vector with coordinate \(b_{ax}\mathbf1\{X=x,A=a,R=1\}/(p_xe_{ax}\widehat r_{n,ax})\).

The pilot observations are bounded, the number of cells is fixed, and each cell probability \(p_xe_{ax}\) is positive. The second-moment calculation in \(\mathrm{thm:coupled\text{-}measurable\text{-}attainment}\), now applied at the natural stochastic scale, gives
\[
\lVert b\rVert=O_P(m_n^{-1/2})
\tag{41}
\]
locally uniformly over the fixed nuisance fiber; the probability of an empty pilot cell tends to zero under \(m_n\to\infty\). Because \(q_n-q^0=s(h,v)/\sqrt n\), finiteness of the support and boundedness of \(s(h,v)\) give
\[
\max_{a,x}|\sigma^2_{ax}(q_n(h,v))-\sigma^2_{ax}(q^0)|=O(n^{-1/2})
\tag{42}
\]
uniformly on that fiber. The \((0,x)\) and \((1,x)\) columns of \(A_{\mathrm{map}}=[-F\;F]\) are \(-p_xd_x\) and \(p_xd_x\), respectively. Consequently
\[
A_{\mathrm{map}}\operatorname{diag}_{a,x}\left\{
\frac{\sigma^2_{ax}(q^0)}{p_xe_{ax}r_{ax}}
\right\}A_{\mathrm{map}}^\top
=\sum_{x=1}^K\sum_{a=0}^1
\frac{p_x\sigma^2_{ax}(q^0)}{e_{ax}r_{ax}}d_xd_x^\top
=\Sigma(r).
\tag{43}
\]
The fixed lower bounds on \(p_xe_{ax}\widehat r_{n,ax}\), (39)--(43), and \(\widehat r_n\to r^\dagger\) now imply
\[
\lVert S_n-\Sigma(r^\dagger)\rVert
=O_P(m_n^{-1}+n^{-1/2})+o_P(1)=o_P(1).
\tag{44}
\]
Thus the covariance remainder uses only \(m_n^{-1}\to0\).

The variance identity (7) in \(\mathrm{thm:coupled\text{-}measurable\text{-}attainment}\), \(q^0_{ax\ell}\ge\eta\), and the two distinct support points available for every \(L\ge2\) give a positive cell-variance floor. Together with \(\operatorname{rank}(F)=d\), equation (43) gives a common positive lower eigenvalue for \(S_n\) with probability tending to one. The summands in (38) are uniformly bounded because \(|Y|\le C_Y\), the pilot means are clipped to the convex hull of \(\mathcal Y\), \(\widehat r_{n,ax}\ge\rho\), and \(p_xe_{ax}>0\). After normalization by \(S_n^{-1/2}\), the sum of third Euclidean absolute moments of the \(n\) scaled summands is \(O(n^{-1/2})\). Applying \(\mathrm{lem:bentkus\text{-}convex\text{-}set}\) and then (44) proves
\[
\sqrt n\{\widetilde\theta_n-\theta_p(q_n(h,v))\}
\Rightarrow N(0,\Sigma(r^\dagger))
\tag{45}
\]
locally uniformly over \(h\in\mathcal H\) and \(\lVert v\rVert\le C_v\).

The revalidated \(\mathrm{thm:local\text{-}chart\text{-}membership}\) gives \(\sqrt n\,\theta_p(q_n(h,v))=h\), including when \(L=2\). Hence the input \(\sqrt n\,\widetilde\theta_n\) to the terminal rule converges to \(N(h,\Sigma(r^\dagger))\). The coupled theorem chooses its rectangle count \(J_{s_n}\) with \(J_{s_n}/\sqrt n\to0\), so the convex-set error summed over its finitely many rectangles is \(o(1)\). Its state-mesh, Gaussian-tail, quadrature, rounding, covariance, and finite-game errors also vanish. Since \(r^\dagger\) minimizes \(L\),
\[
L(r^\dagger)=V^*.
\]
Therefore
\[
\sup_{h\in\mathcal H,\,\lVert v\rVert\le C_v}
E_{P_{p,q_n(h,v)},\mathcal A^*}
[\ell_{\mathrm{reg}}(h,J_n(\mathcal A^*))]
\le V^*+o(1).
\tag{46}
\]
The loss is bounded by \(2\sup_{h\in\mathcal H}\lVert h\rVert_\infty\), so conditional errors converging in probability may be integrated. Taking \(\limsup_{n\to\infty}\) in (46) for fixed \(C_v\), and then \(C_v\to\infty\), gives
\[
\mathcal R_+(\mathcal A^*)\le V^*.
\tag{47}
\]
The established \(\mathrm{thm:local\text{-}minimax\text{-}lower\text{-}bound}\) gives \(\mathcal R_-(\mathcal A)\ge V^*\) for every \(\mathcal A\in\mathfrak A_{\mathrm{det}}\). Since \(\mathcal R_-(\mathcal A^*)\le\mathcal R_+(\mathcal A^*)\), (47) proves that the lower and upper deterministic-limit local minimax envelopes both equal \(V^*\), completing the proof.
